\chapter{Appendix C: \texttt{cloudkube} README}
\label{appendix:cloudkube-readme}




\section*{\texttt{cloudkube} Developer Documentation}

\subsection*{Initial Setup}

The initial setup for \texttt{cloudkube} is straightforward and needs to be done only once. Since \texttt{cloudkube} runs on Python (\texttt{>= 3.11}), it is recommended to use virtual environments to avoid dependency conflicts and keep your global environment clean. Virtual environments can be created using \textbf{conda} (recommended), \textbf{pipenv}, or \textbf{python venv}.

\begin{enumerate}[leftmargin=4pt]
    \item \textbf{Create and Activate a Virtual Environment:}
          Use your preferred method (e.g., conda, pipenv, or venv) to create and activate a virtual environment.

    \item \textbf{Install \texttt{cloudkube}:}
          Run the following command to install \texttt{cloudkube} via \texttt{pip}:
          \begin{lstlisting}
pip install cloudkube
          \end{lstlisting}

    \item \textbf{Verify Installation:}
          Check if \texttt{cloudkube} was installed successfully:
          \begin{lstlisting}
pip list | grep cloudkube
          \end{lstlisting}

    \item \textbf{Install AWS CLI:}
          Follow the official AWS documentation to install the AWS CLI: 
          \href{https://docs.aws.amazon.com/cli/latest/userguide/getting-started-install.html}{\color{blue}\underline{Getting Started with AWS CLI}}.
          Configure your cloud credentials:
          \begin{lstlisting}
aws configure
          \end{lstlisting}
          For additional details, refer to the 
          \href{https://docs.aws.amazon.com/cli/latest/userguide/cli-chap-welcome.html}{\color{blue}\underline{AWS CLI Welcome Guide}}.

    \item \textbf{Install Terraform:}
          \begin{itemize}[leftmargin=4pt]
              \item \textbf{macOS (Recommended):}
                    Install Terraform using Homebrew:
                    \begin{lstlisting}
brew tap hashicorp/tap
brew install hashicorp/tap/terraform
brew update
brew upgrade hashicorp/tap/terraform
                    \end{lstlisting}

              \item \textbf{Other Platforms:}
                    Terraform can also be installed on other platformrs by following the \href{https://developer.hashicorp.com/terraform/tutorials/aws-get-started/install-cli}{\color{blue}\underline{Terraform Developer Documentation}}.
          \end{itemize}

    \item \textbf{Verify Terraform Installation:}
          \begin{lstlisting}
terraform -version
          \end{lstlisting}
\end{enumerate}

Once all dependencies are installed, you are ready to start using \texttt{cloudkube}!

\subsection*{\texttt{cloudkube} Workflow}

\subsubsection*{1. Initialize \texttt{cloudkube}}
To start working with cloudkube, run `cloudkube init`. This will bring up and interactive wizard, which will bring you through the provisioning of a cluster. 
\begin{lstlisting}
cloudkube init
\end{lstlisting}
This starts an interactive wizard to guide you through the provisioning of a Kubernetes cluster.

\subsubsection*{2. Configure Cluster Configuration}
\begin{enumerate}
    \item In the wizard, select \textbf{Option 2} to configure the cluster.
    \item Provide the following information when prompted:
          \begin{itemize}
              \item Name of the cluster
              \item Region
              \item Compute architecture
              \item Optional add-ons (e.g., persistent volumes, ingress controllers)
          \end{itemize}
    \item The configuration details will be saved in a file named \texttt{config.json} in the same directory where \texttt{cloudkube} is run.
\end{enumerate}

\subsubsection*{3. Spin Up the Cluster}
\begin{enumerate}
    \item Select \textbf{Option 3} in the interactive wizard.
    \item \texttt{cloudkube} performs the following actions:
          \begin{itemize}
              \item \texttt{terraform init}: Installs all required Terraform dependencies and sets up the environment.
              \item \texttt{terraform plan}: Creates a resource plan and seeks user approval before proceeding.
              \item \texttt{terraform apply}: Provisions the required infrastructure, including:
                    \begin{itemize}
                        \item An EKS-compliant VPC
                        \item Relevant permissions and security groups
                    \end{itemize}
              \item Automatically configures the local environment to communicate with the cluster by populating \texttt{\textasciitilde/.kube/config}.
          \end{itemize}
\end{enumerate}

\subsubsection*{4. Verify Cluster Creation}
\begin{enumerate}
    \item Check the cluster creation via the \textbf{AWS Console}.
    \item Use the following \texttt{kubectl} commands for further verification:
          \begin{itemize}
              \item List all pods:
                    \begin{lstlisting}
kubectl get pods -A
                    \end{lstlisting}
              \item List all nodes:
                    \begin{lstlisting}
kubectl get nodes
                    \end{lstlisting}
          \end{itemize}
\end{enumerate}

\subsubsection*{5. Apply Kubernetes Manifests}
To deploy your application:
\begin{enumerate}
    \item Ensure your Kubernetes YAML files are stored in the \texttt{./kubernetes} directory.
    \item Apply the manifests:
          \begin{lstlisting}
kubectl apply -f ./kubernetes
          \end{lstlisting}
\end{enumerate}

\subsubsection*{6. Teardown and Resource Deprovisioning}
Once you are done testing your application:
\vspace{-1em}
\begin{enumerate}
    \item Remove your Kubernetes deployment:
          \begin{lstlisting}
kubectl delete -f ./kubernetes
          \end{lstlisting}
    \item Run the following command and select \textbf{Option 4} in the wizard to:
          \begin{itemize}
              \item Deprovision all resources
              \item Clean up your local Kubernetes configuration file (\texttt{\textasciitilde/.kube/config}):
                    \begin{lstlisting}
cloudkube init
                    \end{lstlisting}
          \end{itemize}
\end{enumerate}

\subsection*{Notes}
\begin{itemize}
    \item As of version \texttt{v1.0}, \texttt{cloudkube} only supports \textbf{AWS Cloud}. If support for other cloud providers is added in the future, their respective setup steps should be followed.
    \item Ensure that the EKS cluster meets all networking and security requirements specified in the AWS documentation.
\end{itemize}

